%% ============================================================
%% Revision Notes (2026-03-29, Yuyang):
%% 1. Core Principle: Removed all 6 mathematical formulas. Condensed
%%    Definition + Basic Framework + Theoretical Perspective into a
%%    single concise subsection (~15 lines of prose).
%% 2. Applications: Merged sub-paragraphs (Prompting/RAG/Reasoning/Memory)
%%    into one paragraph per trustworthiness domain (Safety, Privacy,
%%    Fairness, Factuality). Removed \textbf sub-headings.
%% 3. Discussion: Removed from this file. Shortened version moved to
%%    sections/11-challenges_directions.tex with red source marker.
%% ============================================================

\subsection{Context Engineering}
\label{sec:context_engineering}

\begin{figure}[htbp]
    \centering
    \resizebox{\textwidth}{!}{%
        \begin{tikzpicture}
        % --- 1. Define leaf nodes (papers, citations added) ---
        \node[leaf] (safety_papers) {
            Kumar et al.~\citep{kumar2024certifying}, Cao et al.~\citep{cao2024defending}, Zhou et al.~\citep{zhou2024defending}, SemanticSmooth~\citep{ji2024defending}, Wei et al.~\citep{wei2023jailbreak}, Hong et al.~\citep{hong2024so}, Zhong et al.~\citep{zhong2023poisoning}, SafetyAnalyst~\citep{li2025safetyanalyst}, SAFFRON~\citep{qiu2025saffron}.
        };
        \node[leaf, below=0.3cm of safety_papers] (privacy_papers) {
            Wu et al.~\citep{wu2024quantifying}, Chaudhari et al.~\citep{chaudhari2024phantom}, PoisonedRAG~\citep{zou2025poisonedrag}, InferDPT~\citep{tong2025inferdpt}.
        };
        \node[leaf, below=0.3cm of privacy_papers] (fairness_papers) {
            Dwivedi et al.~\citep{dwivedi2023breaking}, Huang et al.~\citep{huang2024bias}, FACTER~\citep{fayyazi2025facter}, SPICA~\citep{chen2024spica}, FairRAG~\citep{shrestha2024fairrag}, Dhingra et al.~\citep{dhingra2023queer}, Kamruzzaman et al.~\citep{kamruzzaman2024prompting}.
        };
        \node[leaf, below=0.3cm of fairness_papers] (factuality_papers) {
            Dhuliawala et al.~\citep{dhuliawala2024chain}, Press et al.~\citep{press2023measuring}, Think Twice~\citep{li2024think}, Feng et al.~\citep{feng2024teaching}, Zhou et al.~\citep{zhou2023context}, REPLUG~\citep{shi2024replug}, SELF-RAG~\citep{asai2024self}, DELL~\citep{wan2024dell}, Yu et al.~\citep{yu2023generate}, MetaRAG~\citep{zhou2024metacognitive}, Chen et al.~\citep{chen2024improving}, Wang et al.~\citep{wang2023augmenting}.
        };

        % --- 2. Define main sub-nodes (Trustworthiness dimensions) ---
        \node[sub, left=0.5cm of safety_papers] (safety) {Safety};
        \node[sub, left=0.5cm of privacy_papers] (privacy) {Privacy};
        \node[sub, left=0.5cm of fairness_papers] (fairness) {Fairness};
        \node[sub, left=0.5cm of factuality_papers] (factuality) {Factuality \&  Faithfulness};

        % --- 3. Define root node ---
        \node[main, left=2.5cm of {$(safety.west)!0.5!(factuality.west)$}] (ce) {Context Engineering (§\ref{sec:context_engineering})};

        % --- 4. Draw connection lines ---
        % a. Draw a main trunk from the root node to a branching point
        \coordinate (fork) at ($(ce.east) + (1.25cm,0)$);
        \draw[line] (ce.east) -- (fork);

        % b. Connect from main branching point to each sub-node
        \draw[line] (fork) |- (safety.west);
        \draw[line] (fork) |- (privacy.west);
        \draw[line] (fork) |- (fairness.west);
        \draw[line] (fork) |- (factuality.west);

        % c. Connect main sub-nodes directly to leaf nodes
        \draw[line] (safety.east) -- (safety_papers.west);
        \draw[line] (privacy.east) -- (privacy_papers.west);
        \draw[line] (fairness.east) -- (fairness_papers.west);
        \draw[line] (factuality.east) -- (factuality_papers.west);

        \end{tikzpicture}%
    }
    \caption{A taxonomy of Context Engineering for enhancing LLM trustworthiness.}
    \label{fig:ce_taxonomy}
\end{figure}

Within the landscape of inference-time methods for trustworthy LLMs, context engineering (CE) has emerged as a central paradigm. Instead of modifying model parameters, CE emphasizes the deliberate design and orchestration of the information presented to the model at inference time. This approach is particularly important for trustworthiness: by shaping the context, one can systematically influence factuality, safety, fairness, etc, without modifying model's parameters.

\subsubsection{Core Principle}

Context engineering is the systematic discipline of designing, optimizing, and managing the information payload provided to LLMs at inference time. Unlike early prompt engineering, which treated the context as a flat string, CE models the context as a structured composition of four components: directive instructions encompassing rules and constraints; knowledge inputs such as retrieved documents or prior dialogue; interaction mechanisms including tools, APIs, or system affordances; and the live user query specifying the task at hand. An assembly function integrates these heterogeneous elements into the final sequence consumed by the model.

The CE framework rests on three foundational pillars that together form a pipeline. Context retrieval and generation is responsible for sourcing task-relevant material through prompt design, retrieval from external knowledge bases, and adaptive assembly strategies. Context processing transforms raw material into model-ready input through compression, refinement, and multimodal or structured integration. Context management ensures coherence and efficiency across stages by dynamically organizing, prioritizing, and updating contextual elements under system constraints such as the maximum context length. Together, these pillars support higher-level methods such as Retrieval-Augmented Generation (RAG), persistent memory, tool-augmented reasoning, and multi-agent collaboration.


\subsubsection{Applications}

\paragraph{Safety}

Context engineering promotes safety by controlling the informational boundary within which a model operates, encoding explicit constraints at the prompt level, grounding generation in trusted external sources, and embedding verification routines into the reasoning process. Prompt engineering offers a lightweight first layer of defense, ranging from certifiable defenses with robustness guarantees against adversarial prompts~\citep{kumar2024certifying} and robust alignment prompting~\citep{cao2024defending} to baseline defenses such as paraphrasing, in-context refusals, and perplexity checks~\citep{zhou2024defending}, as well as smoothing-based methods like SemanticSmooth that guard against both token- and prompt-level attacks~\citep{ji2024defending}. Retrieval-augmented generation further enhances safety by grounding responses in verifiable external knowledge, reducing hallucinations and lowering the risk of unsafe misinformation~\citep{wei2023jailbreak, hong2024so}, though vulnerabilities such as corpus poisoning attacks that insert adversarial passages into knowledge bases remain a concern~\citep{zhong2023poisoning}. Reasoning-oriented approaches complement these strategies by enabling models to regulate their own outputs during inference through structured step-by-step verification~\citep{li2025safetyanalyst} and safety-aware scoring that resists jailbreaks at scale~\citep{qiu2025saffron}.

\paragraph{Privacy}

Context engineering for privacy focuses on controlling what enters the model's context. Instead of modifying parameters, privacy is preserved by assembling the context such that sensitive elements are masked, filtered, or abstracted while retaining task-relevant information. Recent work shows that CE can both introduce and mitigate privacy risks. Prompts themselves may leak sensitive training data through property and membership inference attacks~\citep{wu2024quantifying}, and RAG pipelines are vulnerable to poisoned documents that can exfiltrate private passages or bias model outputs~\citep{chaudhari2024phantom}. On the defense side, certified privacy-preserving mechanisms for RAG offer robustness against adversarial manipulation~\citep{zou2025poisonedrag}, while frameworks for selective retrieval and privacy-aware generation control personal data exposure~\citep{tong2025inferdpt}. Together, these studies underscore that privacy risks often stem from how external context is engineered into models, and effective mitigation requires careful control, sanitization, and certification of injected information.

\paragraph{Fairness}

Bias in LLMs often arises from skewed training data and emergent model behavior, and context engineering provides a practical means of mitigation without retraining. Carefully designed prompts have been shown to reduce bias, from employing in-context learning to mitigate gender stereotypes~\citep{dwivedi2023breaking} to frameworks for testing and reducing social bias in code generation through iterative, feedback-driven prompting~\citep{huang2024bias} and embedding explicit fairness constraints in prompts~\citep{fayyazi2025facter}. Retrieval-augmented methods enhance fairness by grounding outputs in curated, balanced evidence, with fairness-aware retrieval mechanisms prioritizing diverse and representative sources to reduce stereotype reinforcement~\citep{chen2024spica, shrestha2024fairrag}. A third line of work integrates structured reasoning through chain-of-thought~\citep{dhingra2023queer} or deliberative reasoning~\citep{kamruzzaman2024prompting} to enforce fairness constraints, decomposing decisions into interpretable steps that enable bias detection and correction during inference.

\paragraph{Factuality}

A core challenge in trustworthy AI is mitigating factual errors and hallucinations, and context engineering offers several complementary strategies. Prompt engineering improves factual consistency through structured instructions, including self-detection prompts that require models to evaluate multiple candidate answers before finalizing a response~\citep{dhuliawala2024chain, press2023measuring, li2024think}, abstention-oriented prompting that encourages models to decline uncertain answers in multilingual settings~\citep{feng2024teaching}, and opinion-based or counterfactual demonstrations that reduce reliance on parametric knowledge~\citep{zhou2023context}. Retrieval-augmented generation enhances factual grounding by supplementing LLMs with external documents, with frameworks enabling adaptive and reflective retrieval such as SELF-RAG~\citep{asai2024self}, integration of retrieval with proxy explanations for misinformation detection~\citep{wan2024dell, yu2023generate, shi2024replug}, and metacognitive loops that detect insufficient or conflicting knowledge and adjust retrieval accordingly~\citep{zhou2024metacognitive}. Memory-based methods further improve factuality by allowing models to retain and update knowledge across interactions through external memory buffers where user feedback persists across sessions~\citep{chen2024improving} and modular memory systems combining episodic traces and semantic caches~\citep{wang2023augmenting}.
